\par\Needspace{6\baselineskip}\medskip
\hypertarget{native-124-article091}{}\pdfbookmark[2]{91.}{article-091}
\begin{center}91.\end{center}
Aequatio 12, in qua, eo de quo agimus casu, manifesto signum superius adoptare oportet, per introductionem quantitatis $\xi$ obtinet formam
\[
m=(l+x)^{\frac12}+\dfrac{(l+x)^{\frac32}}{\frac34-\frac{9}{10}(x-\xi)}\mathpunct{\cdot}
\]
Statuendo itaque $\BareRoot{(l+x)}=\dfrac{m}{y}\mathpunct{\raisebox{.45ex}{,}}$ atque
\[
\mathrm{14)}\qquad \dfrac{mm}{\frac56+l+\xi}=h,
\]
omnibus rite reductis prodit
\[
\mathrm{15)}\qquad h=\dfrac{(y-1)yy}{y+\frac13}\mbox{\HistoricalDisplayNoteMark}.
\]
\noindent\HistoricalDisplayNoteText{Handschriftliche Bemerkung:]
\[
y-1=\dfrac{l+x}{\frac34-\frac{9}{10}(x-\xi)}
\]
\[
y=1+\frac{10}{9}h-\frac{110}{81}hh+\frac{2410}{729}h^3-\text{etc.}
\]
\[
yy=1+\frac{20}{9}h-\frac{40}{27}hh+\frac{2620}{729}h^3-\text{etc.}
\]
\[
\dfrac{1}{yy}=1-\frac{20}{9}h+\frac{520}{81}hh-\frac{5140}{243}h^3+\text{etc.}
\]
}Quodsi itaque $h$ tamquam quantitatem cognitam spectare licet, $y$ inde per aequationem cubicam determinabitur, ac dein erit
\[
\mathrm{16)}\qquad x=\dfrac{mm}{yy}-l.
\]
Iam etiamsi $h$ implicet quantitatem adhuc incognitam $\xi$, in approximatione \SourceLeaf{125}prima eam negligere atque pro $h$ accipere licebit $\dfrac{mm}{\frac56+l}\mathpunct{\raisebox{.45ex}{,}}$ quoniam certe in eo de quo agimus casu $\xi$ semper est quantitas valde parva. Hinc per aequationes 15, 16 elicientur $y$ et $x$; ex $x$ per tabulam III habebitur $\xi$, cuius adiumento per formulam 14 eruetur valor correctus ipsius $h$, cum quo calculus idem repetitus valores correctos ipsarum $y$, $x$ dabit: plerumque hi tam parum a praecedentibus different, ut $\xi$ iterum e tabula III desumta haud diversa sit a valore primo; alioquin calculum denuo repetere oporteret, donec nullam amplius mutationem patiatur. Simulac quantitas $x$ inventa erit, habebitur $g$ per formulam $\sin\tfrac12 g^2=x$.

Haec praecepta referuntur ad casum primum, ubi $\cos f$ positivus est; in casu altero ubi negativus est statuimus $\BareRoot{(L-x)}=\dfrac{M}{Y}$ atque
\[
\mathrm{14}^{\ast}\mathrm{)}\qquad \dfrac{MM}{L-\frac56-\xi}=H,
\]
unde aequatio $12^{\ast}$ rite reducta transit in hanc
\[
\mathrm{15}^{\ast}\mathrm{)}\qquad H=\dfrac{(Y+1)YY}{Y-\frac13}\mathpunct{\cdot}
\]
Per hanc itaque aequationem cubicam determinare licet $Y$ ex $H$, unde rursus $x$ derivabitur per aequationem
\[
\mathrm{16}^{\ast}\mathrm{)}\qquad x=L-\dfrac{MM}{YY}\mathpunct{\cdot}
\]
In approximatione prima pro $H$ accipietur valor $\dfrac{MM}{L-\frac56}\mathpunct{\raisebox{.45ex}{;}}$ cum valore ipsius $x$ inde per aequationes $15^{\ast}$, $16^{\ast}$ derivato desumetur $\xi$ ex tabula III; hinc per formulam $14^{\ast}$ habebitur valor correctus ipsius $H$, cum quo calculus eodem modo repetetur. Tandem ex $x$ angulus $g$ eodem modo determinabitur ut in casu primo.
